% PROJECT SUMMARY -- upload as a separate one-page document.
\section*{Project Summary}

\subsection*{Overview}
Ice-sheet science increasingly combines numerical models, ensemble data
assimilation, observational archives, and scalable computing, but researchers
still spend substantial effort installing software, translating data,
rewriting scheduler workflows, and reconstructing computational context. This
three-year CSSI Elements project will advance \textbf{CryoStack} from a
promising developer-operated stack into secure, reusable, and measurable
cyberinfrastructure for reproducible cryosphere computing. Existing code
establishes feasibility: password/GitHub/ORCID identity, role-based
administration and a Control Center, persistent workspaces and experiment
histories, a modularized CryoLauncher interface, ICESEE data assimilation,
LIVIST radar/borehole temperature exploration, a connector/relay path to
Slurm and (for CryoLauncher's ISSM workflow on AWS) to an institutional MATLAB
license, a manifest-driven Frozen Legacies historical-radar catalog,
and a bring-your-own-account AWS path on which CryoLauncher's ISSM and Icepack
workflows and a single-process ICESEE Lorenz-96 example have run end to end. It is not yet a finished community
service: extension contracts, production-safe remote execution, scientifically
qualified cloud and HPC runs, archival provenance, operations, governance, and
adoption evidence remain incomplete.

The project has three objectives. \textbf{O1} will deliver versioned
application/model/dataset extension kits, automated CryoBench qualification, a
production PISM coupling, and a second Frozen Legacies collection.
\textbf{O2} will harden an authentication-bound, allow-listed execution plane
for local systems, Georgia Tech PACE, and project- or user-managed AWS
accounts, with administrator-controlled computing allocations. \textbf{O3} will create durable experiment/provenance services, a
connected historical-data-to-model workflow, measurable operations, user
evaluation, training, and contributor governance. Each extension will reuse
shared identity, validation, execution, provenance, and deployment services.

\subsection*{Intellectual Merit}
CryoStack will enable reproducible experiments in which common observation
bundles and assimilation configurations can be evaluated across PISM, ISSM,
Icepack, and idealized models, supporting studies of ice-sheet state,
uncertain basal and rheological parameters, and structural uncertainty. The CI
innovation is a clean, versioned separation of scientific and dataset
adapters, execution backends, and gateway/operations services, validated by
PISM and a second historical-data collection, neither of which shaped the
interfaces. CryoBench
will gate every extension on schema, numerical, scientific, security, and
data-quality tests, including \textbf{cross-backend scientific equivalence}:
whether a workflow remains scientifically consistent, within stated
tolerances, across local, institutional-HPC, and cloud execution. A hardened, allow-listed control plane will
replace the current broad-shell connector for securely bridging the gateway
to protected institutional resources, and every run will produce a versioned
experiment/result/provenance record spanning configuration, execution,
backend, and application-specific diagnostics. PISM and one additional
historical collection are the committed new integrations; the same contracts
create a bounded path for future community models, data, and analysis tools
without rebuilding the surrounding infrastructure.

\subsection*{Broader Impacts}
Browser-first, containerized workflows will lower infrastructure barriers for
students and researchers without dedicated research-software support. Three
annual workshops, reusable tutorials, office hours, and paired contribution
reviews will train participants in data assimilation, reproducibility, HPC,
and scientific software engineering, prioritizing early-career researchers and
under-resourced institutions. Open releases and archival DOIs, accessible
documentation, transparent governance, privacy-preserving metrics, and a
five-year operations plan will establish a credible path from working
prototypes to community-maintained cyberinfrastructure. The primary alignment
is CISE/OAC; the domain alignment is GEO/OPP.
